\chapterAndLabel{Modifiability Checker}{modifiability-checker}

\newcommand{\UnsupportedOperationException}{\<Un\-sup\-port\-ed\-Op\-er\-a\-tion\-Ex\-cep\-tion>}

In the Java Collections Framework, an \emph{optional} method is one that a class may or may not support.
If a method is unsupported, calling it throws \UnsupportedOperationException.
The Modifiability Checker verifies that optional methods related to modifying collections are only
called if they are supported, thus eliminating the possibility of an
\UnsupportedOperationException\  at run time.
(It is possible to write a checker that warns about other unsupported
methods, but the Modifiability Checker focuses on those issued when modifying collections.)

An example of an optional method is
\sunjavadoc{java.base/java/util/Collection.html\#add(E)}{add}.  Some
subclasses of \<Collection> support \<add>, and others (such as the
result of calling
\sunjavadoc{java.base/java/util/Collections.html\#unmodifiableCollection(java.util.Collection)}{Collections.unmodifiableCollection})
do not support \<add>.

If the Modifiability Checker issues no warning, then your program is
guaranteed not to throw \UnsupportedOperationException\ due to invoking an
optional method (such as
\sunjavadoc{java.base/java/util/Collection.html\#add(E)}{add} or
\sunjavadoc{java.base/java/util/Collection.html\#remove(java.lang.Object)}{remove})
on an unmodifiable collection.
(For brevity, this chapter generally uses
just ``collection'' to mean ``collection, map, or iterator''.)

In addition to preventing UOEs, the Modifiability Checker ensures that collections
that the user intends to be unmodifiable have an unmodifiable run-time implementation.


To run the Modifiability Checker:

\begin{Verbatim}
javac -processor org.checkerframework.checker.modifiability.ModifiabilityChecker MyFile.java
\end{Verbatim}

\sectionAndLabel{Unmodifiable collections}{modifiability-unmodifiable-collections}

An \emph{unmodifiable collection} is one that throws
\UnsupportedOperationException\ if a mutating method is called.
The JDK provides many such collections, for example, the results of:
\begin{itemize}
\item
  \sunjavadoc{java.base/java/util/List.html\#of(E...)}{\code{List.of()}},
  \sunjavadoc{java.base/java/util/List.html\#copyOf(java.util.Collection)}{\code{List.copyOf()}}
\item
  \sunjavadoc{java.base/java/util/Set.html\#of(E...)}{\code{Set.of()}},
  \sunjavadoc{java.base/java/util/Set.html\#copyOf(java.util.Collection)}{\code{Set.copyOf()}}
\item
  \sunjavadoc{java.base/java/util/Map.html\#of(K,V,K,V)}{\code{Map.of()}},
  \sunjavadoc{java.base/java/util/Map.html\#copyOf(java.util.Map)}{\code{Map.copyOf()}}
\item
  \sunjavadoc{java/util/Collections.html\#emptyList()}{\code{Collections.emptyList()}},
  \sunjavadoc{java/util/Collections.html\#emptySet()}{\code{Collections.emptySet()}},
  \sunjavadoc{java/util/Collections.html\#emptyMap()}{\code{Collections.emptyMap()}}
\item
  \sunjavadoc{java.base/java/util/Collections.html\#unmodifiableList(java.util.List)}{\code{Collections.unmodifiableList()}},
  \sunjavadoc{java.base/java/util/Collections.html\#unmodifiableSet(java.util.Set)}{\code{Collections.unmodifiableSet()}},
  \sunjavadoc{java.base/java/util/Collections.html\#unmodifiableMap(java.util.Map)}{\code{Collections.unmodifiableMap()}}
\end{itemize}

A \emph{modifiable collection} supports calling at least one optional mutating method. In general,
classes support groups of methods which correspond to four different ways a collection can be modified.
A collection may support modifying a collection through one or more of the capabilities
\textbf{Grow}, \textbf{Sequential Grow}, \textbf{Shrink}, and \textbf{Replace}.

\begin{itemize}
    \item \textbf{Grow}: The collection can grow by adding new elements
      (\<add>, \<addAll>, \<offer>, \ldots) at arbitrary locations.
       Not to be confused with sequential grow.
    \item \textbf{Sequential Grow}: The sequenced collection can grow by adding new
      elements at the beginning or end (\<addFirst>, \<addLast>,
      \<offerFirst>, \<offerLast>, \<push>, \<putFirst>,
      \<putLast>, \ldots).

      The Sequential Grow capability is needed for collections, such as
      \sunjavadoc{java.base/java/util/TreeSet.html}{\code{TreeSet}}, that
      support grow operations (like \<add> and \<addAll>) but do not
      support sequential grow operations (like \<addFirst> and
      \<addLast>).  The declaration of \<TreeSet> is annotated as \<@Growable @SeqUngrowable>.

      ``Sequential Grow'' is shortened to ``SeqGrow'' in the rest of the manual.
    \item \textbf{Shrink}: The collection can shrink by removing
      existing elements (\<remove>, \<removeAll>, \<clear>, \<retainAll>, \<poll>, \<pop>, \ldots).
    \item \textbf{Replace}: The collection can replace existing
      elements with new values at specific locations/keys (\<List.set>, \<Map.replace>, \<Map.replaceAll>,
       \<ListIterator.set>,  \<Collections.sort>, \<Map.Entry.setValue>, \ldots).
\end{itemize}

\noindent
Some collections can also be modified through their iterator; see Section~\ref{modifiability-iterator-behavior}.

The following table summarizes the main method families that may be called
for each capability.  It is not a list of every overload.  Read-only
methods such as \<get>, \<contains>, \<size>, \<isEmpty>, \<iterator>, and
\<entrySet> do not require a modifiability capability.

\noindent
\begin{tabular}{lp{0.75\linewidth}}
Capability required & Method families \\
\<@Growable> &
  \<Collection.add>, \<addAll>, \<List.add>,
  \<List.addAll>, \<Queue.add>, \<Queue.offer>, \<BlockingQueue.put>,
  timed \<offer>, \<TransferQueue.transfer>, \<tryTransfer>,
  \<Map.putIfAbsent>, \<Map.computeIfAbsent>, \<ListIterator.add>,
  legacy \<Vector> insertion methods, and \<Collections.addAll> \\
\<@SeqGrowable> &
  \<SequencedCollection.addFirst>, \<addLast>, \<Deque.addFirst>,
  \<addLast>, \<offerFirst>, \<offerLast>, \<push>,
  \<BlockingDeque.putFirst>, \<putLast>, timed \<offerFirst>,
  timed \<offerLast>, \<SequencedMap.putFirst>, \<putLast> \\
\<@Shrinkable> &
  \<Collection.remove>, \<removeAll>, \<retainAll>,
  \<removeIf>, \<clear>, \<List.remove>, \<Queue.remove>, \<poll>,
  \<BlockingQueue.take>, timed \<poll>, \<Deque.pop>, \<removeFirst>,
  \<removeLast>, \<pollFirst>, \<pollLast>, occurrence-removal methods,
  \<Map.remove>, conditional \<remove>, \<Map.clear>,
  \<NavigableMap.pollFirstEntry>, \<pollLastEntry>, \<Iterator.remove>,
  and legacy \<Vector> removal methods \\
\<@Replaceable> &
  \<List.set>, \<List.replaceAll>,
  \<List.sort>, \<Map.replace>, \<Map.replaceAll>,
  \<Map.Entry.setValue>, \<ListIterator.set>, and mutating
  \<Collections> algorithms such as \<sort>, \<reverse>, \<shuffle>,
  \<swap>, \<fill>, \<copy>, \<rotate>, and \<replaceAll>, plus
  legacy \<Vector> replacement methods \\
Combined capabilities &
  \<Map.put> and \<putAll> require \<@Growable @Replaceable>;
  \<Map.computeIfPresent> requires \<@Shrinkable @Replaceable>;
  \<Map.compute> and \<merge> require
  \<@Growable @Shrinkable @Replaceable>;
  \<BlockingQueue.drainTo> requires a \<@Shrinkable> receiver and a
  \<@Growable> target collection \\
\end{tabular}

An overriding method must preserve positive modifiability receiver
requirements exactly.  For example, if a superclass method has a
\<@Growable>, \<@SeqGrowable>, \<@Shrinkable>, or \<@Replaceable> receiver,
then the overriding method must require that same positive receiver
capability; it may not drop or weaken the requirement.

\sectionAndLabel{Modifiability annotations}{modifiability-annotations}

The Modifiability Checker uses \textbf{multiple independent 4-element
hierarchies}, one per capability.  Every collection type carries exactly
one qualifier from each hierarchy.  The checker also uses an additional
2-element hierarchy to express whether a collection's iterator preserves the
collection's capabilities:  whether \<remove()> can be called on the iterator,
and, for a \<ListIterator>, whether \<add()> and \<set()> can be called.

\begin{figure}
\begin{smaller}
  \centering
  \begin{minipage}{0.85\linewidth}
\begin{verbatim}
  Grow hierarchy           SeqGrow hierarchy             Shrink hierarchy              Replace hierarchy             Iterator hierarchy

    @MaybeGrowable          @MaybeSeqGrowable         @MaybeShrinkable               @MaybeReplaceable           @MaybeIteratorPolyMod
      /      \                 /       \                    /       \                    /       \                          |
 @Growable @Ungrowable   @SeqGrowable @SeqUngrowable  @Shrinkable @Unshrinkable   @Replaceable @Unreplaceable               |
      \      /                 \       /                    \       /                    \       /                          |
    @BottomGrowable        @BottomSeqGrowable          @BottomShrinkable              @BottomReplaceable            @IteratorPolyMod
\end{verbatim}
  \end{minipage}
\end{smaller}
  \caption{Independent modifiability qualifier hierarchies.
    Each annotated type carries one qualifier from each hierarchy.}
  \label{fig-modifiability-lattice}
\end{figure}

Figure~\ref{fig-modifiability-lattice} shows the hierarchies.
Each hierarchy is type-checked individually.
An assignment \code{T x = y} is valid only if the compile-time type of \code{y} is a
subtype of (or equal to) \code{T} in every hierarchy simultaneously.  For
example, \code{\emph{@Growable} @MaybeSeqGrowable @MaybeShrinkable @MaybeReplaceable List} is a
subtype of \code{\emph{@MaybeGrowable} @MaybeSeqGrowable @MaybeShrinkable @MaybeReplaceable List}
(because \<@Growable> $\le$ \<@MaybeGrowable>), but it is \emph{not} a
subtype of \code{\emph{@Ungrowable} @MaybeSeqGrowable @MaybeShrinkable @MaybeReplaceable List}
(because \<@Growable> and \<@Ungrowable> are incomparable siblings in the
Grow hierarchy).

\subsectionAndLabel{Alias annotations}{modifiability-aliases}

For convenience, the Modifiability Checker provides five alias
annotations.  They are not part of any hierarchy; the checker expands
each one into its constituent qualifiers when it appears in source code.

\begin{description}

\item[\refqualclass{checker/modifiability/qual}{Modifiable}]
Alias for \code{@Growable @SeqGrowable @Shrinkable @Replaceable} (but see
Section \ref{modifiability-types-lacking-methods} for exceptions).  Calling
any mutating operation on this collection will not throw
\UnsupportedOperationException.

\item[\refqualclass{checker/modifiability/qual}{Unmodifiable}]
Alias for \code{@Ungrowable @SeqUngrowable @Unshrinkable @Unreplaceable} (but see
Section~\ref{modifiability-types-lacking-methods} for exceptions).  The
collection is definitely unmodifiable: calling any mutating operation
always throws \UnsupportedOperationException.  The checker issues a warning
at every invocation of a grow, seq-grow, shrink, or replace operation.

\item[\refqualclass{checker/modifiability/qual}{MaybeModifiable}]
Alias for \code{@MaybeGrowable @MaybeSeqGrowable @MaybeShrinkable @MaybeReplaceable}.
The checker cannot determine the modifiability of the collection.
The checker conservatively issues a warning at every invocation of a grow,
seq-grow, shrink, or replace operation.

\item[\refqualclass{checker/modifiability/qual}{UnmodifiableParam}]
Syntactic sugar for \refqualclass{checker/modifiability/qual}{MaybeModifiable}.
It may only be written within a formal parameter type to indicate that the
method does not modify the parameter.  This annotation (and
\refqualclass{checker/modifiability/qual}{MaybeModifiable}) enable the
programmer to distinguish between a data structure's capability to perform
a modification and the programmer's intent about whether a modification
should occur.
\refqualclass{checker/modifiability/qual}{UnmodifiableParam}
may appear nested within a parameter type, such as in
\code{List<@UnmodifiableParam List<String>\relax>}.

\refqualclass{checker/modifiability/qual}{UnmodifiableParam} represents
\emph{reference unmodifiability}:  no modification may happen through the
reference, but an alias may modify the value.  By contrast,
\refqualclass{checker/modifiability/qual}{Unmodifiable} is object
unmodifiability:  no alias may modify the value.

\item[\refqualclass{checker/modifiability/qual}{PolyModifiable}]
Alias for \code{@PolyGrowable @PolySeqGrowable @PolyShrinkable @PolyReplaceable}
(see Section~\ref{modifiability-polymorphic}).

\end{description}

\noindent
Examples using alias annotations:

\begin{Verbatim}
  @Modifiable List<String> mod = new ArrayList<>();  // all four capabilities
  @Unmodifiable List<String> unmod = List.of("a");   // no capabilities
  @MaybeModifiable List<String> unknown = ...;  // capabilities are unknown at compile time
  void foo(@UnmodifiableParam List<String> list) { ... }  // foo does not modify list
\end{Verbatim}

\subsectionAndLabel{Qualifiers}{modifiability-qualifiers}

\subsubsectionAndLabel{The Grow hierarchy}{modifiability-grow-hierarchy}

\begin{description}
\item[\refqualclass{checker/modifiability/qual}{MaybeGrowable}]
The top qualifier in the Grow hierarchy.
The checker cannot determine whether this collection supports grow
operations.
Calling grow operations such as \code{add}, \code{addAll}, etc.\ on this
collection \emph{may} throw \UnsupportedOperationException.
This is the default qualifier for unannotated types in the
Grow hierarchy.

\item[\refqualclass{checker/modifiability/qual}{Growable}]
Calling grow operations such as \code{add}, \code{addAll}, etc.\ on this
collection \emph{will not} throw \UnsupportedOperationException.

\item[\refqualclass{checker/modifiability/qual}{Ungrowable}]
Calling grow operations such as \code{add}, \code{addAll}, etc.\ on this
collection \emph{will} throw \UnsupportedOperationException.

\item[\refqualclass{checker/modifiability/qual}{BottomGrowable}]
The bottom qualifier in the Grow hierarchy.  Programmers should rarely write it.
\end{description}

\subsubsectionAndLabel{The SeqGrow hierarchy}{modifiability-seqgrow-hierarchy}

\begin{description}
\item[\refqualclass{checker/modifiability/qual}{MaybeSeqGrowable}]
The top qualifier in the SeqGrow hierarchy.
The checker cannot determine whether this collection or map supports sequenced grow
operations.
Calling sequenced grow operations such as \code{addFirst}, \code{addLast},
\code{putFirst}, and \code{putLast} on this collection or map \emph{may}
throw \UnsupportedOperationException.
This is the default qualifier for unannotated types in the
SeqGrow hierarchy.

\item[\refqualclass{checker/modifiability/qual}{SeqGrowable}]
Sequenced grow operations, such as \code{addFirst} and \code{addLast} for
collections or \code{putFirst} and \code{putLast} for maps, \emph{will not}
throw \UnsupportedOperationException.

\item[\refqualclass{checker/modifiability/qual}{SeqUngrowable}]
Sequenced grow operations, such as \code{addFirst} and \code{addLast} for
collections or \code{putFirst} and \code{putLast} for maps, \emph{will}
throw \UnsupportedOperationException.
For example, \<TreeSet> and \<ConcurrentSkipListSet> support ordinary
\<add> but throw \UnsupportedOperationException\ for explicit positional
insertion with \<addFirst> and \<addLast>.
Similarly, \<TreeMap> and \<ConcurrentSkipListMap> support ordinary \<put>
but throw \UnsupportedOperationException\ for explicit positional updates
with \<putFirst> and \<putLast>.

\item[\refqualclass{checker/modifiability/qual}{BottomSeqGrowable}]
The bottom qualifier in the SeqGrow hierarchy.  Programmers should rarely write it.
\end{description}

\subsubsectionAndLabel{The Shrink hierarchy}{modifiability-shrink-hierarchy}

\begin{description}
\item[\refqualclass{checker/modifiability/qual}{MaybeShrinkable}]
The top qualifier in the Shrink hierarchy.
The checker cannot determine whether this collection supports shrink
operations.
Calling shrink operations such as \code{remove}, \code{clear}, etc.\ on
this collection \emph{may} throw \UnsupportedOperationException.
This is the default qualifier for unannotated types in the
Shrink hierarchy.

\item[\refqualclass{checker/modifiability/qual}{Shrinkable}]
Calling shrink operations such as \code{remove}, \code{clear}, etc.\ on
this collection \emph{will not} throw \UnsupportedOperationException.

\item[\refqualclass{checker/modifiability/qual}{Unshrinkable}]
Calling shrink operations such as \code{remove}, \code{clear}, etc.\ on
this collection \emph{will} throw \UnsupportedOperationException.

\item[\refqualclass{checker/modifiability/qual}{BottomShrinkable}]
The bottom qualifier in the Shrink hierarchy.  Programmers should rarely write it.
\end{description}

\subsubsectionAndLabel{The Replace hierarchy}{modifiability-replace-hierarchy}

\begin{description}
\item[\refqualclass{checker/modifiability/qual}{MaybeReplaceable}]
The top qualifier in the Replace hierarchy.
The checker cannot determine whether this collection supports replace
operations.
Calling replace operations such as \code{set}, \code{replaceAll}, etc.\
on this collection \emph{may} throw \UnsupportedOperationException.
This is the default qualifier for unannotated types in the
Replace hierarchy.

\item[\refqualclass{checker/modifiability/qual}{Replaceable}]
Calling replace operations such as \code{set}, \code{replaceAll}, etc.\
on this collection \emph{will not} throw \UnsupportedOperationException.

\item[\refqualclass{checker/modifiability/qual}{Unreplaceable}]
Calling replace operations such as \code{set}, \code{replaceAll}, etc.\
on this collection \emph{will} throw \UnsupportedOperationException.

\item[\refqualclass{checker/modifiability/qual}{BottomReplaceable}]
The bottom qualifier in the Replace hierarchy.  Programmers should rarely write it.
\end{description}

\subsubsectionAndLabel{The Iterator hierarchy}{modifiability-iterator-hierarchy}

\begin{description}
\item[\refqualclass{checker/modifiability/qual}{MaybeIteratorPolyMod}]
The top qualifier.
The checker cannot determine whether this collection's \code{iterator()} is
\<@PolyShrinkable>.
This is the default qualifier for unannotated types.

\item[\refqualclass{checker/modifiability/qual}{IteratorPolyMod}]
This collection's \code{iterator()} is \<@PolyShrinkable>.  That is, if
collection \<c> is \<@Shrinkable>, then \<c.iterator()> is also
\<@Shrinkable>.

\end{description}

It might seem that it is sufficient to simply annotate \<iterator()>
directly, as in

\begin{Verbatim}
class List {
  @MaybeShrinkable Iterator iterator() { ... }
}
class ArrayList {
  @PolyShrinkable Iterator iterator(@PolyShrinkable ArrayList this) { ... }
}
class CopyOnWriteArrayList {
  @Unshrinkable Iterator iterator(@PolyShrinkable ArrayList this) { ... }
}
\end{Verbatim}

That approach would be sound but imprecise.
In practice, many expressions of static type \<List> evaluate to
\<ArrayList>s or other classes with \<@PolyShrinkable> \<iterator()> methods.
The \<@IteratorPolyMod> annotation is similar to writing the hypothetical
annotation \<@RuntimeType("ArrayList")>:  it indicates a property that is
more precise than the declared type.


\subsectionAndLabel{Polymorphic qualifiers}{modifiability-polymorphic}

A polymorphic qualifier (Section~\ref{method-qualifier-polymorphism})
specifies that a method's return type has the
same qualifier (in some hierarchy) as a formal parameter type (possibly
the receiver).

\begin{description}

\item[\refqualclass{checker/modifiability/qual}{PolyGrowable}]
Polymorphic qualifier for the Grow hierarchy.
The return type's growability matches the growability of whichever
formal parameter is annotated with \<@PolyGrowable>.

\item[\refqualclass{checker/modifiability/qual}{PolySeqGrowable}]
Polymorphic qualifier for the SeqGrow hierarchy.
The return type's seq-growability matches the seq-growability of
whichever formal parameter is annotated with \<@PolySeqGrowable>.

\item[\refqualclass{checker/modifiability/qual}{PolyShrinkable}]
Polymorphic qualifier for the Shrink hierarchy.
The return type's shrinkability qualifier matches the shrinkability of
whichever formal parameter is annotated with \<@PolyShrinkable>.
Useful for methods such as \code{Map.keySet()} that preserve shrinkability while
returning an \<@Ungrowable> view.

\item[\refqualclass{checker/modifiability/qual}{PolyReplaceable}]
Polymorphic qualifier for the Replace hierarchy.
The return type's replaceability matches the replaceability of
whichever formal parameter is annotated with \<@PolyReplaceable>.

\item[\refqualclass{checker/modifiability/qual}{PolyIteratorPolyMod}]
Polymorphic qualifier for the Iterator hierarchy.
The return type's iterator-preservation qualifier matches the
iterator-preservation qualifier of whichever formal parameter is annotated
with \<@PolyIteratorPolyMod>.

\item[\refqualclass{checker/modifiability/qual}{PolyModifiable}]
Alias for \code{@PolyGrowable @PolySeqGrowable @PolyShrinkable @PolyReplaceable}.
Preserves all four capabilities simultaneously.
Use on methods such as \code{Collections.synchronizedList} that
do not change modifiability.

\end{description}


\subsectionAndLabel{Modifiability method annotations}{modifiability-method-annotations}

The Modifiability Checker supports method annotations that specify method
behavior.

\begin{description}

\item[\refqualclass{checker/modifiability/qual}{PreservesModifiability}]
Indicates that if the argument to the method is modifiable, then the returned collection is
also modifiable. However, if the argument is unmodifiable, then the returned collection may or may
not be unmodifiable. This is useful for methods that may return their argument, or may return a
new modifiable collection.
For example, if the argument is \<@Growable>, then the returned collection is also
\<@Growable>. But if the argument is \<@Ungrowable>, then the return collection is \<@MaybeGrowable>.
More generally, if the argument is \<@Growable>, \<@Shrinkable>,
\<@Replaceable>, \<@SeqGrowable>, or \<@IteratorPolyMod>, then the Modifiability Checker
treats the return value as having that same capability. If the argument has
any other qualifier, the return type is the top qualifier (\<@MaybeGrowable>,
\<@MaybeSeqGrowable>, \<@MaybeShrinkable>, \<@MaybeReplaceable>, or \<@MaybeIteratorPolyMod>) in
the corresponding hierarchy.

This annotation may only be written on non-void, one-argument method declarations.

\end{description}

\sectionAndLabel{\<@Modifiable> and \<@Unmodifiable> for types without grow, seq-grow, shrink, and/or replace methods}{modifiability-types-lacking-methods}

Ordinarily, the declaration
\<@Modifiable MyClass x;> is equivalent to
\<@Growable @SeqGrowable @Shrinkable @Replaceable MyClass x;>,
and
\<@Unmodifiable MyClass x;> is equivalent to
\<@Ungrowable @SeqUngrowable @Unshrinkable @Unreplaceable MyClass x;>.

However, some types lack methods for one or more mutation
capabilities.  In such cases, the Modifiability Checker weakens the
structurally unavailable component to the corresponding top qualifier.
For example, \<Iterator> has
\<remove> but not \<add> or \<set>, so \<@Modifiable Iterator> expands to
\<@MaybeGrowable @Shrinkable @MaybeReplaceable Iterator>.


\subsectionAndLabel{Alternative design: use negative capability rather than top qualifier}{modifiability-modifiable-must-be-top}

An alternative design would make \<Modifiable Iterator> expand to
\<@Ungrowable @Shrinkable @Unreplaceable Iterator>.
That would be incorrect, because not every \<Iterator> is \<@Ungrowable> and \<@Unreplaceable>.
For example, the static type \<Iterator> does not declare \<add> or \<set>
(they are structurally unavailable), but its subtype
\<ListIterator> and does declare those methods.
Therefore \<@Modifiable Iterator> should not mean that every
iterator is incapable of grow and replace; it means only that those capabilities
are not known from the static type \<Iterator>.

Similarly, \<Queue> does not declare replacement methods, but a variable
with declared type
\<Queue> may contain a \<LinkedList>, which does support
replacement operations.

Using the top qualifier for structurally unavailable methods ensures that
code like this type-checks:

\begin{Verbatim}
@Replaceable LinkedList list = ...;
Queue q = list;
@Replaceable LinkedList list2 = (LinkedList) q; // still @Replaceable
\end{Verbatim}


\subsectionAndLabel{Expansion rules for \<@Modifiable> and \<@Unmodifiable>}{modifiability-modifiable-expansion-rules}

The following rules define this weakening.  The first qualifier in each
capability row is used for \<@Modifiable>; the second is used for
\<@Unmodifiable>.  If the weakening condition holds, both aliases use the
corresponding \<@Maybe*> qualifier instead.



\noindent
\begin{tabular}{l|l|l}
\multicolumn{3}{c}{Expansion for
\<@Modifiable>/\<@Unmodifiable>. ``$\backslash$'' means set difference.  A
type includes its subtypes.} \\ \hline
& Usual expansion to \<@>(\<Un>)\<*> & Expands to \<@Maybe*> \\ \hline
Grow
& \<ListIterator>
& \<Map.Entry> or (\<Iterator>$\backslash$\<ListIterator>) \\
SeqGrow
& \<SequencedCollection>, \<SequencedMap>, or \<Deque>
& not \<SequencedCollection>, \<SequencedMap>, or \<Deque> \\
Shrink
& not \<Map.Entry>
& \<Map.Entry> \\
Replace
& \<LinkedList> or \<ListIterator>
& exact \<Collection>, any \<Set>, (\<Queue>$\backslash$\<LinkedList>), \\
&&  or (\<Iterator>$\backslash$\<ListIterator>) \\
\end{tabular}

These rules describe alias expansion, not the full specification of every JDK
type.  The JDK annotations may still give a more precise capability for a
particular declaration.  For instance, \<SortedSet> is a
\<SequencedCollection> and \<SortedMap> is a \<SequencedMap>, so their
seq-grow component is structurally relevant and is not weakened to
\<@MaybeSeqGrowable>.  Their JDK declarations are annotated
\<@SeqUngrowable>, because their positional insertion methods such as
\<addFirst>, \<addLast>, \<putFirst>, and \<putLast> always throw
\UnsupportedOperationException.


\subsectionAndLabel{Example expansions for \<@Modifiable> and \<@Unmodifiable>}{modifiability-modifiable-expansion-examples}

The following are representative canonical expansions.

\noindent
\begin{smaller}
\setlength{\tabcolsep}{.33\tabcolsep}
\begin{tabular}{rlcrrrrl}
\<@Modifiable> & \<Collection> & = & \<@Growable> & \<@MaybeSeqGrowable> & \<@Shrinkable> & \<@MaybeReplaceable> & \<Collection> \\
\<@Unmodifiable> & \<Collection> & = & \<@Ungrowable> & \<@MaybeSeqGrowable> & \<@Unshrinkable> & \<@MaybeReplaceable> & \<Collection> \\
\<@Modifiable> & \<Queue>    & = & \<@Growable> & \<@MaybeSeqGrowable> & \<@Shrinkable> & \<@MaybeReplaceable> & \<Queue> \\
\<@Unmodifiable> & \<Queue>    & = & \<@Ungrowable> & \<@MaybeSeqGrowable> & \<@Unshrinkable> & \<@MaybeReplaceable> & \<Queue> \\
\<@Modifiable> & \<Deque>    & = & \<@Growable> & \<@SeqGrowable> & \<@Shrinkable> & \<@MaybeReplaceable> & \<Deque> \\
\<@Unmodifiable> & \<Deque>    & = & \<@Ungrowable> & \<@SeqUngrowable> & \<@Unshrinkable> & \<@MaybeReplaceable> & \<Deque> \\
\<@Modifiable> & \<Map>      & = & \<@Growable> & \<@MaybeSeqGrowable> & \<@Shrinkable> & \<@Replaceable> & \<Map> \\
\<@Unmodifiable> & \<Map>      & = & \<@Ungrowable> & \<@MaybeSeqGrowable> & \<@Unshrinkable> & \<@Unreplaceable> & \<Map> \\
\<@Modifiable> & \<SequencedMap> & = & \<@Growable> & \<@SeqGrowable> & \<@Shrinkable> & \<@Replaceable> & \<SequencedMap> \\
\<@Unmodifiable> & \<SequencedMap> & = & \<@Ungrowable> & \<@SeqUngrowable> & \<@Unshrinkable> & \<@Unreplaceable> & \<SequencedMap> \\
\<@Modifiable> & \<Set>      & = & \<@Growable> & \<@MaybeSeqGrowable> & \<@Shrinkable> & \<@MaybeReplaceable> & \<Set> \\
\<@Unmodifiable> & \<Set>      & = & \<@Ungrowable> & \<@MaybeSeqGrowable> & \<@Unshrinkable> & \<@MaybeReplaceable> & \<Set> \\
\<@Modifiable> & \<Iterator> & = & \<@MaybeGrowable> & \<@MaybeSeqGrowable> & \<@Shrinkable> & \<@MaybeReplaceable> & \<Iterator> \\
\<@Unmodifiable> & \<Iterator> & = & \<@MaybeGrowable> & \<@MaybeSeqGrowable> & \<@Unshrinkable> & \<@MaybeReplaceable> & \<Iterator> \\
\<@Modifiable> & \<ListIterator> & = & \<@Growable> & \<@MaybeSeqGrowable> & \<@Shrinkable> & \<@Replaceable> & \<ListIterator> \\
\<@Unmodifiable> & \<ListIterator> & = & \<@Ungrowable> & \<@MaybeSeqGrowable> & \<@Unshrinkable> & \<@Unreplaceable> & \<ListIterator> \\
\<Map.@Modifiable> & \<Entry> & = & \<@MaybeGrowable> & \<@MaybeSeqGrowable> & \<@MaybeShrinkable> & \<@Replaceable> & \<Map.Entry> \\
\<Map.@Unmodifiable> & \<Entry> & = & \<@MaybeGrowable> & \<@MaybeSeqGrowable> & \<@MaybeShrinkable> & \<@Unreplaceable> & \<Map.Entry> \\
\end{tabular}
\end{smaller}


\sectionAndLabel{Behavior of the \<iterator()> method}{modifiability-iterator-behavior}

By default, the signature of \<Collection.iterator()> is

\begin{Verbatim}
@MaybeShrinkable Iterator iterator()
\end{Verbatim}

However, for some subclasses of \<Collection>, the signature of
\<iterator()> is

\begin{Verbatim}
@PolyShrinkable Iterator iterator(@PolyShrinkable MySubclassOfCollection this)
\end{Verbatim}

A collection uses the second signature if the collection is
\<@IteratorPolyMod>.  The collection uses the first signature if the
collection is \<@MaybeIteratorPolyMod>.

Here are examples of its use:

\begin{Verbatim}
  @IteratorPolyMod @Shrinkable List<String> list = ...;
  @Shrinkable Iterator<String> it = list.iterator();
  it.remove();  // OK

  @Unshrinkable List<String> list2 = ...;
  @Unshrinkable Iterator<String> it2 = list2.iterator();
  it2.remove();  // Error: it2 is @Unshrinkable

  @Shrinkable List<String> list3 = ...;
  @MaybeShrinkable Iterator<String> it3 = list3.iterator();
  it3.remove();  // Error: list3 does not have @IteratorPolyMod
\end{Verbatim}

% \sectionAndLabel{Unverified collection implementation annotations}{modifiability-class-unverified}
%
% The Modifiability Checker verifies uses of collection modifiability annotations,
% but it does not verify collection implementations.
%
% When a user-defined collection-like class (a \<Collection>, \<Map>,
% \<Iterator>, or subtype) explicitly writes a concrete modifiability qualifier
% on the class declaration or on a constructor, the checker issues a
% \<class.unverified> warning.  The warning means that the
% annotation is a trusted specification of the custom implementation, not a fact
% verified from the implementation's method bodies.  After examining the
% implementation, suppress the warning with
% \<@SuppressWarnings("class.unverified")>.
%
% For example, the class-level \<@Growable> annotation below triggers
% \<class.unverified>. It is up to the programmer to verify that
% the implementation of \<TestMod> is not \@Growable and remove it from the class declaration.
% \begin{Verbatim}
% public @Growable class TestMod extends ArrayList<String> {
%
%   @Override
%   public boolean add(@Growable TestMod this, String s) {
%     throw new UnsupportedOperationException();
%   }
% }
% \end{Verbatim}



\sectionAndLabel{Examples of Modifiability Checker warnings}{modifiability-examples}

\begin{Verbatim}
import java.util.*;
import org.checkerframework.checker.modifiability.qual.*;

class Demo {

  void basicUsage() {
    // @Modifiable is an alias for @Growable @SeqGrowable @Shrinkable @Replaceable
    @Modifiable List<String> mod = new ArrayList<>();
    mod.add("a");     // OK: @Growable guarantees add() works
    mod.addFirst("z"); // OK: @SeqGrowable guarantees addFirst() works
    mod.remove("a");  // OK: @Shrinkable guarantees remove() works
    mod.set(0, "b");  // OK: @Replaceable guarantees set() works

    // @Unmodifiable is an alias for @Ungrowable @SeqUngrowable @Unshrinkable @Unreplaceable
    @Unmodifiable List<String> unmod = List.of("a", "b");
    unmod.add("c");   // Error: add() requires @Growable, got @Ungrowable
    unmod.addFirst("d"); // Error: addFirst() requires @SeqGrowable, got @SeqUngrowable
    unmod.get(0);     // OK: read-only access needs no capability
  }

  void fineGrainedPermissions(@Growable List<String> g,
                              @Shrinkable List<String> s,
                              @Replaceable List<String> r) {
    // Growable list allows adding elements
    g.add("a");       // OK
    g.remove("a");    // Error: remove() requires @Shrinkable
    g.set(0, "b");    // Error: set() requires @Replaceable

    // Shrinkable list allows removing elements
    s.remove("a");    // OK
    s.add("a");       // Error: add() requires @Growable

    // Replaceable list allows updating elements
    r.set(0, "b");    // OK
    r.add("b");       // Error: add() requires @Growable
  }

  void combinedPermissions(@Growable @Replaceable Map<String, String> map) {
    // Map.put requires both Grow and Replace capabilities
    map.put("key", "value");  // OK

    // Map.remove requires Shrink capability
    map.remove("key");        // Error: @MaybeShrinkable (default) !<: @Shrinkable
  }

  // @PolyModifiable preserves all four capabilities
  @PolyModifiable List<String> wrap(@PolyModifiable List<String> list) {
    return list;
  }

  void testPoly(@Modifiable List<String> mod, @Growable @Shrinkable List<String> gs) {
    @Modifiable List<String> m = wrap(mod);            // OK
    @Growable @Shrinkable List<String> gs2 = wrap(gs); // OK
    // :: error: [assignment]
    @Modifiable List<String> bad = wrap(gs);  // Error: gs has @MaybeSeqGrowable and @MaybeReplaceable
  }
}
\end{Verbatim}

% LocalWords:  Modifiability modifiability copyOf emptyList emptySet gsr
% LocalWords:  emptyMap unmodifiableList unmodifiableSet unmodifiableMap
% LocalWords:  toArray addAll removeAll retainAll replaceAll setValue Un
% LocalWords:  MaybeModifiable UnmodifiableParam Growable MaybeGrowable MaybeShrinkable
% LocalWords:  MaybeReplaceable MaybeSeqGrowable growable seq ListIterator
% LocalWords:  PolyModifiable PolySeqGrowable addFirst synchronizedList Un
% LocalWords:  PolyShrinkable PolyGrowable PolyReplaceable SeqGrow Deque
% LocalWords:  addLast offerFirst offerLast putFirst putLast unreplaceable seq-growability
% LocalWords:  SequencedMap SortedMap unmodifiableCollection entrySet
% LocalWords:  SeqGrowable SeqUngrowable BottomSeqGrowable Growability shrinkability replacability
% LocalWords:  growability Shrinkable Replaceable Ungrowable Unshrinkable
% LocalWords:  Unreplaceable BottomGrowable BottomShrinkable BottomReplaceable
% LocalWords:  MaybeIteratorPolyMod IteratorPolyMod PolyIteratorPolyMod PreservesModifiability
% LocalWords:  BlockingQueue TransferQueue tryTransfer putIfAbsent putAll
% LocalWords:  SequencedCollection BlockingDeque removeIf removeFirst
% LocalWords:  removeLast pollFirst pollLast NavigableMap pollFirstEntry
% LocalWords:  pollLastEntry computeIfPresent drainTo unmodifiability
% LocalWords:  seqgrow ConcurrentSkipListSet ConcurrentSkipListMap keySet
% LocalWords:  RuntimeType MyClass SortedSet
